% Registrations: the six registration documents, what each covers, when they were filed, what the
% authors had seen of each study's data before filing, what each status label means, how the images
% were made, and the corrections to earlier reports. Numbers come only from macros in
% generated/numbers.tex.
\label{sec:suppA-registrations}

A registration is a dated record of a study's hypotheses and analysis plan. It has nothing to do with copyright registration. This section lists the six registration documents of the five studies and what each covers. It says what we had seen of each study's data when we filed them, defines the status labels that every result in the paper carries, states how the images were made, and corrects statements made in earlier reports of these studies.

\subsection*{The six documents}

All six documents were filed together on 8 October 2026 on the Open Science Framework (OSF), an online archive for study plans, as one registration in the project at \url{https://osf.io/u3d69/overview?view_only=be795b275c3548428d74457df31a4159}. Each study has a short registration that answers the nine questions of the form used by AsPredicted, another registration service: whether data had been collected, the hypotheses, the measures, the conditions, the analyses, the exclusions, the sample size, anything else to record, and a title. Study~1 also has a long registration, which sets out its full analysis plan, the legal reasoning behind its predictions and the rules for its text measures. Where Study~1's short registration is silent, the long one governs. Supplementary Table~\ref{tab:suppA-registrations} lists the six documents, the study each covers and what each registers. Copies of the six files are shared with the data (main text, Data availability).

% The six registration documents: the study each covers, when its data were collected, and what
% it registers. Design facts from the registrations; numbers only through generated/numbers.tex.
% See supplement/tables/README_suppA.md.
\begin{table}[tbp]
\centering
\small
\caption{\textbf{The six registration documents.} All six were filed together on the Open Science Framework (OSF) on 8 October 2026, after every study's data had been collected. Studies are listed in the paper's order; the order in which they were run was 2, 4, 1, 3, 5. \emph{Confirmatory}: the small set of tests that a registration names as the study's confirmatory tests, judged together after Holm's correction. Bounds are in rating points on 1-to-5 scales.}
\label{tab:suppA-registrations}
\begin{tabular}{@{}>{\raggedright\arraybackslash}p{3.0cm}>{\raggedright\arraybackslash}p{3.1cm}>{\raggedright\arraybackslash}p{\dimexpr\textwidth-6.1cm-4\tabcolsep\relax}@{}}
\toprule
Document & Study, data collection and analyzed sample & What it registers\\
\midrule
Study~1, short registration & Study~1; 25 February to 3 March 2025; $N$~=~\SOneN & H1, the AI label lowers both creativity and effort; H2, the drop is larger on effort; H3, the drop is the same for images with and without a human figure; H4, under the AI label ratings vary more between people and less within a person. Four confirmatory tests (two for H1, one each for H2 and H4); H3 decided by its own rule against $\pm$\SOneHThreeBound. Five word-meaning tests on the explanations as an exploratory set, and the two text checks.\\
\addlinespace
Study~1, long registration & Study~1 & The full plan, which governs where the short registration is silent: why the predictions matter for the law (section~2), design, exclusions, every model and decision rule, the robustness checks, the word lists and text analyses (section~9), which decisions were taken with the estimates already known (section~8), and the materials (section~10).\\
\addlinespace
Study~2 & Study~2; 23 November to 5 December 2024; $N$~=~\STwoN & H1, ratings of the second work rise from None to Slight to Dramatic; H2, the label's effect, averaged over levels, lies inside $\pm$\SourceBound{} on each measure; H3, whether the label's effect shrinks from None to Dramatic (no direction predicted). Twelve confirmatory tests, four per measure; ordinal checks; the two text checks on creativity explanations.\\
\addlinespace
Study~3 & Study~3; 4 to 9 March 2025; $N$~=~\SThreeN & The label, the stated time and their interaction on creativity and on effort (six confirmatory $F$~tests); an equivalence test of the Human-minus-AI gap at 100 hours against $\pm$\SourceBound; a rule deciding whether time closes the gap, the gap persists, or the result is inconclusive; ordinal checks; the two text checks on both ratings.\\
\addlinespace
Study~4 & Study~4; 27 January to 2 February 2025; $N$~=~\SFourN & H1 and H2 as in Study~2; H3, whether the hours stated for the original change ratings of the second work, alone or with the level (no direction predicted). Fifteen confirmatory tests, five per measure. Secondary: the label's effect at each level, its change with stated hours, an equivalence test of 100 hours against under 1 hour, and a refit without Human-label participants who wrote that a machine made the work; the two text checks on creativity explanations.\\
\addlinespace
Study~5 & Study~5; 21 October to 2 November 2025; $N$~=~\SFiveN & The effect of stated editing hours, and its interaction with the label, on each measure (six confirmatory $F$~tests); an equivalence test of editing time against $\pm$\SFiveEditTimeBound; the label's effect at each level against $\pm$\SourceBound, nine tests each judged at .05 without correction; ordinal checks; a second analysis with every item rescaled from 0 to 1; the two text checks on creativity explanations.\\
\bottomrule
\end{tabular}
\end{table}


Every document was filed after data collection had ended in every study. The study numbers follow the argument of the paper, not the order in which the studies were run. Study~2 ran first, from 23 November to 5 December 2024, then Study~4, Study~1 and Study~3 in early 2025, and Study~5 from 21 October to 2 November 2025. In this paper, `registered' therefore means written into one of these documents. It does not mean fixed before the data were collected, and no study in this paper was registered before its data were collected.

For Studies 2 to 5, each registration was written before the analysis code reported here was run on the real data. Until the filing, that code was run only blind: every model and table was produced with the condition labels shuffled at random across participants, so that the code could be checked while no number described the real conditions. The runs with the real labels followed on 8 October 2026, after the filing. One exception concerns the written explanations. In Study~5's blind runs, and in one early blind run of Study~3's, the ratings were not shuffled, so the text analyses that relate the words to the rating alone used the real ratings. Their outputs were set aside, and no estimate from them is reported anywhere. Study~1's registration is different: it was written after its analyses had been run, which makes it a postregistration.

\subsection*{What we had seen before filing}

The first question of each registration asks whether data had been collected, and each answer lists what we had already done with them. The account below follows those answers.

\paragraph*{Study~1.}
In April 2025 we compared the two label groups' mean ratings with Welch $t$~tests, which compare two means without assuming that the groups vary equally, and with analyses of variance by label and kind of image, on the \SOneN{} participants that the registered exclusions keep. The first author showed these results at a lab meeting on 18 April 2025. A public talk by the three authors on 26 September 2025 showed the same $t$~tests, with effect sizes, for both ratings. Both drops on the originals, and the larger drop on effort, had therefore been seen before Study~1 was registered. In November 2025 we ran topic models and sentiment scores on the explanations. By April 2026 a language model had coded the explanations with the label, and for sentiment the rating, in its instructions, and a manuscript draft reported the means, effect sizes and robustness checks along with those codes. None of that text work is used here. Before 7 September 2026, other exploratory work had predicted the ratings from the words. In September and October 2026 we ran the registered code itself, and a code review on 8 September 2026 summarized its results for the first author. Section~8 of the long registration names only three elements fixed before those analyses: the model for each rating, the sequence of exclusions and the convention for effect sizes. It lists the decisions taken with the estimates in view. They include the joint model that tests H2, H3's bound and decision table, the set of tests corrected together, and the quantity that decides H4; the rule for fits that did not converge was added in October 2026. Study~1's H1 and H2 are therefore registered tests of patterns already seen, not predictions made in advance.

\paragraph*{Study~2.}
By February 2025 we had fitted an analysis of variance to the \STwoN{} participants that the registered exclusions keep. The CogSci 2025 paper reported it as its first study,\cite{kamper2025originalsin} and the September 2025 talk showed it with plots of every condition. Ratings rose steeply with the level of modification. The label made a difference only to creativity, and only averaged over levels, not at any single level. No interaction was significant. In April 2025 we ran keyword, topic and sentiment analyses of the explanations. After that, only counts and a few pooled answer frequencies were read from the data. Study~2's registered test of the label's effect on creativity therefore re-tests an effect already known.

\paragraph*{Study~3.}
In April 2025 we ran a mixed analysis of variance, one with factors that vary both between and within participants, on the finished records, including some that the registered exclusions later removed, and we saw its mean for each label and stated time. The September 2025 talk reported effects of the label and of stated time on both ratings, and a change in the label's effect with time on effort but not on creativity; it called the label's effect significant at some stated times and not at others. No analysis of the explanations had been run. We wrote the hypotheses and the plan knowing these results.

\paragraph*{Study~4.}
In February 2025 we fitted mixed analyses of variance to these data. The CogSci 2025 paper reported them as its second study, and the September 2025 talk reported them as well. They showed a large effect of modification, no significant main effect of the label, an effect of the hours stated for the original on creativity only, and an effect of the label on essence change that varied with the level of modification. We found no analysis of the explanations. For the registration we read counts, screening answers, image slots and answer codes, and summarized no rating by condition. We wrote the hypotheses knowing the earlier results.

\paragraph*{Study~5.}
Study~5 was run after the CogSci 2025 paper and the talk, and neither reports it. A November 2025 manuscript draft reported a sample size that the registered exclusions do not reproduce, with age and gender, and no rating result. In 2026, coding sessions in which an AI assistant worked for the first author displayed the exclusion counts and cell sizes (7 September), one participant's three ratings and explanations for one trial (22 September), and a pooled count of answers in the extra category of two six-point items (6 October). No rating had been summarized by condition, and no model of the ratings had been fitted. Study~5 is therefore the only study registered before any of its ratings had been averaged or compared by label. Both of its equivalence bounds were fixed before any Study~5 model was fitted; the bound on editing time is read from Study~3's registered result.

\subsection*{What each status label means}

Every result in the paper carries one of the statuses defined below. Supplementary Table~\ref{tab:supp-verdicts} gives the status and verdict of each registered test.

\emph{Registered confirmatory}: one of the small set of tests that a registration names as the study's confirmatory tests. It is judged by the registration's own rule, usually after Holm's correction.\cite{holm1979simple} That correction raises each $p$-value in the set, so that the chance of calling any result in the set significant by mistake stays at .05. A $p$-value is the chance of a result at least this far from no effect if there were no effect. Words such as `supported', `holds', `not supported' and `inconclusive' are verdicts under that rule.

\emph{Registered secondary}: an analysis that a registration names but places outside the confirmatory set, such as the robustness checks, the two text checks run in every study, and, in Studies 4 and 5, the label's effect at each level of modification. These are judged without correction unless the registration says otherwise. Study~1's H3 is a registered hypothesis decided outside the confirmatory set by its own rule, a significance test read together with an equivalence test; Supplementary Table~\ref{tab:supp-verdicts} labels its status `own rule'. An equivalence test asks whether an effect is shown to lie inside a stated bound. It passes when the effect's 90\% confidence interval, the range of values that the data do not rule out at the 90\% level, lies wholly inside the bound, which shows that the effect is small, not that it is zero.\cite{schuirmann1987comparison,lakens2017equivalence}

\emph{Exploratory}: an analysis that a registration names as exploratory, or that no registration names. It carries no verdict. Where a registration fixes a correction for an exploratory set, as Study~1's does for its five word-meaning tests, the corrected results are descriptive readings, not verdicts. The analyses that pool Studies 2, 4 and 5, the image check, the breakdown of Study~5's interaction, and the analyses of talk of origin and of how far explanations moved were specified on 9 and 10 October 2026, after every registered result was known; all are exploratory.

\emph{Descriptive}: a count, a mean, a share or a fitting status, with no test. A \emph{registered description} is a descriptive result that a registration asks us to report.

Two terms describe how a model fitted. A model that did not converge stopped without settling on a solution; it yields no $p$-value or verdict, and its check is reported as not completed. A singular fit converged but estimated a variance or a correlation at the edge of its possible range, a sign that the data cannot support that many quantities; it keeps its numbers and is reported as singular.

\subsection*{How the images were made}

The ten original images were generated with OpenAI's DALL·E, used through ChatGPT. The long Study~1 registration (section~10) and the CogSci 2025 paper both name DALL·E as the generator, and both are right.\cite{kamper2025originalsin} Eight of the ten original files carry a provenance record: a statement of where the image came from, which the service that made the image wrote into the file and signed digitally, so that a later change to the record would show. These records follow C2PA, an industry standard for such records set by the Coalition for Content Provenance and Authenticity. Each of the eight records was signed by OpenAI. Each names DALL·E as the software that created the image, and its two entries were written by OpenAI's programming interface and by ChatGPT. Each also classes the image as made by a trained algorithm, the term these records use for content that an AI system generated. They are dated 6 to 8 November 2024, before the first study began. The other two original files carry no provenance record, so for those two the generator rests on our account. The image check (Supplementary Note~4) extracted the images that the surveys displayed, and the original files match them: eight pixel for pixel, among them the two files without a record, and two with only slight differences in pixel values.

We wrote the prompts ourselves; no file records them.

We made the Slight and Dramatic versions of each original by hand, in Adobe Photoshop, an image-editing program that we operated manually. One Dramatic version was instead exported from Dzine, an online image editor that we used only as a manual editor, and the Dzine logo on the exported file was removed before the surveys showed it. No generative or other AI editing tool changed any image. Of the edited files that we kept, seven of the ten Slight files record the program that saved them, and each names Adobe Photoshop; the other three, and all ten Dramatic files, record none. For those files the program rests on our account, as the generator does for the two originals without a provenance record. The None second work was the original's own image file, shown again unchanged. Every original, in both label groups and in every study, was therefore generated by an AI system, so the Human label, which said that Alex made the originals with design and sketching tools, was a cover story.

\subsection*{Corrections to earlier reports}

The statements below, made in earlier reports of these studies, are wrong, and this paper does not repeat them. Most come from the CogSci 2025 paper,\cite{kamper2025originalsin} our proceedings paper that reported Studies 2 and 4 as its Study~1 and Study~2.

\paragraph*{The source of the prompts.}
The CogSci 2025 paper says that the originals were generated from established prompts in a validated set developed by Orwig and colleagues.\cite{orwig2026using} That is wrong: we wrote the prompts ourselves. The long Study~1 registration (section~10) gives the correct account, that our approach drew on that study but that the prompts were our own.

\paragraph*{Power analysis.}
The CogSci 2025 paper says that the sample size of its first study, this paper's Study~2, was set by a power analysis run in the program G*Power, and that its second study, this paper's Study~4, followed the same criteria. A power analysis calculates how many participants a study needs in order to detect an effect of a given size. No power analysis was run for any of the five studies. The planned sample sizes were the quotas set in each survey, a number of participants for each label group or, in Study~5, for each combination of label and editing time; each study's section gives them. The simulations of what Study~1's design could detect (Supplementary Note~1) were computed with its registered analyses, after its data had been collected, and did not set its sample size.

\paragraph*{Pre-registration.}
The CogSci 2025 paper described its first study, this paper's Study~2, as pre-registered, and it linked to a page that it said held all the pre-registrations. Neither Study~2 nor Study~4 was registered with any registry before its data were collected. Their only planning document was an OSF preregistration template filled in with Google Docs. It was begun on 5 November 2024, before Study~2 ran, last edited on 29 January 2025, while Study~4's data were being collected, and never submitted. Its plan also differed from the registered analyses: it planned analyses of variance rather than mixed models, had no rule excluding participants who said they had not taken part seriously, and for Study~4 planned a measure of perceived effort that the survey does not contain. The Study~4 registration says that the CogSci 2025 submission called Study~4 pre-registered too; the published paper makes the claim of its first study.

\paragraph*{The item wording in Study~2.}
The CogSci 2025 paper quotes the essence-change and creativity items of its first study, this paper's Study~2, with definitions, and says that every item used a five-point scale. In the Study~2 survey only the transformativeness item carried a definition, and two transformativeness items had six answer points (Section~\ref{sec:suppB-study2}).

\paragraph*{The Slight edits.}
The CogSci 2025 paper described the Slight versions as adjustments of hue, coloration and contrast, and the registrations of Studies 2, 4 and 5 describe them as adjustments of color and contrast. Of the ten Slight versions, \XSImageCheckSlightReframedN{} were also cropped or resized (Supplementary Note~4).

\paragraph*{Study~4's sample.}
The CogSci 2025 paper gives the sample of its second study, this paper's Study~4, as \SFourNBeforeValidity{} participants. That is the count left after the first four registered exclusion steps. The export held \SFourNRecords{} records, and those steps removed \SFourNPreviewRemoved{} survey test runs, \SFourNUnfinishedRemoved{} unfinished records and \SFourNNotSeriousRemoved{} participants who did not answer that they had taken part seriously; the fourth step, a wrong answer to an arithmetic question, removed nobody. The registered exclusions then remove \SFourNInvalidRemoved{} more records that are not valid completions: one participant rated all ten originals rather than six, and another's record names one original in two of its six image slots. That leaves \SFourN. The CogSci 2025 paper's own degrees of freedom for this study fit the smaller number, which is also the number of participants that the Study~4 registration gives for the February 2025 analyses, all but one of whom are in the registered sample.

\paragraph*{Viewing time.}
The CogSci 2025 paper says that each pair of an original and its second work stayed on screen for at least 8 seconds. In the surveys of Studies 2 and 4, as in Study~5, each such pair stayed on screen for at least 10 seconds; the 8-second minimum applied to an original shown alone.

\paragraph*{The law.}
The CogSci 2025 paper described the US Copyright Office's current policy as treating AI assistance as an `original sin', and it read its ratings as judgments of the qualities that copyright protection requires. This paper does neither. The Office refuses copyright registration for material whose expressive elements an AI system determined. It does register a person's own creative selection or arrangement of that material, or a modification of it creative enough to meet copyright's standard of originality, with the AI-generated material left out of the claim.\cite{usco2023guidance,usco2025copyrightability} In this paper, `original sin' names a hypothesis about lay judgment, that an AI origin passes to whatever is built from it; it is not a legal rule. The ratings are lay judgments of creativity, effort, transformativeness and essence change, not judgments of authorship or of whether a work can be protected.

\paragraph*{Reading the CogSci 2025 paper beside this one.}
Four further points are not errors, but a reader comparing the two papers needs them. First, the names differ: the CogSci 2025 paper's Study~1 is this paper's Study~2, and its Study~2 is this paper's Study~4. Its creator-type factor is the label, and its effort-level factor is the number of hours stated for the original. Second, in Study~2 both papers analyze the same \STwoN{} participants, left after the same two exclusions: \STwoNUnfinishedRemoved{} unfinished records and \STwoNNotSeriousRemoved{} participants who did not answer that they had taken part seriously. Study~4's sample is corrected above. Third, the analyses differ: the CogSci 2025 paper fitted analyses of variance to each participant's average ratings, while the registered analyses fit mixed models to every rating, with a baseline for each participant and each original, and correct their confirmatory tests together. Fourth, one registered result differs: the CogSci 2025 paper reported that the hours stated for the original changed creativity ratings of the second work. Study~4's registered model does not reproduce that effect ($F$(\SFourCreatSourceTimeDfNum,~\SFourCreatSourceTimeDfDen)~=~\SFourCreatSourceTimeF, $p$~\SFourCreatSourceTimeP{} before correction; Holm $p$~\SFourCreatSourceTimePHolm). Some exploratory comparisons also differ, and the sections on Studies 2 and 4 set out the two analyses side by side.

One number in a registration is also not used. Study~1's registration gives the effect sizes that the September 2025 talk showed. An effect size here is Cohen's $d$, the difference between the label groups' mean ratings divided by the spread of participants' mean ratings. The talk's value for creativity matches the analysis reported here (\SOneDCreat), but its value for effort matches no table that the analysis writes, which gives \SOneDEff. The paper uses the values that the analysis gives.
